\chapter{Vergleich von Testausgabeformaten}\label{Vergleich Testausgabefomate}
Um das Tool \texttt{testAuditor} zu entwickeln, welches framework-agnostisch eingesetzt werden kann, muss \texttt{testAuditor} die Ausgabe eines Testlaufes unabhängig vom verwendeten Framework parsen können. Die Anforderungen an ein Testausgabeformat wurden bereits in \chapref{Anforderungen Testausgabeformat} untersucht.\\
Standardmäßig geben die meisten Testframeworks das Ergebnis eines Testlaufs in einem unstrukturierten Format über die Konsole aus. Allerdings bieten fast alle etablierten Testframeworks auch eine Auswahl an Dateiformaten an, über welche die Testausgabe ebenfalls präsentiert werden kann.\\
Für diese Dateiformate wurden über die Jahre einige framework-übergreifende Ausgabeformate entwickelt, die sich in unterschiedlicher Ausprägung durchsetzen konnten.\\
Es gilt nun, ein geeignetes Testausgabeformat zu finden, welches von \texttt{testAuditor} als Quelle für Testergebnisse genutzt werden kann. 

\section{Methodik}
Aufgrund der Vielzahl von verschiedenen Testframeworks und der Vielfältigkeit ihrer Ausgabeformate können im Rahmen dieser Bachelorarbeit nur die verbreitetsten Formate untersucht werden.\\
Zur Untersuchung der praktischen Umsetzung einzelner Formate und zur exemplarischen Erfassung ihrer Verbreitung wird ein stichprobenartiger Datensatz betrachtet. Dabei gilt zu beachten, dass bei diesem Datensatz nur eine kleine Menge ausgewählter Programmiersprachen umfasst. Außerdem beschränken sich die untersuchten Testframeworks auf die am weitesten verbreiteten in ihrer jeweiligen Sprache, um die Etablierung von Testausgabeformaten möglichst praxisnahe darstellen zu können. Außerdem bieten länger etablierte Testframeworks erfahrungsgemäß eine breitere Auswahl von Testausgabeformaten, als weniger verbreitete oder neue Testframeworks.\\
Bei der Untersuchung von Plugins für zusätzliche Testausgabeformate wurden ausschließlich Plugins berücksichtigt, welche über den hier benutzten package manager \texttt{Nix} bereitgestellt werden. Würden externe Plugins, welche zumeist aus unabhängig gepflegten Repositorien stammen, ebenfalls betrachtet werden, würde die Zahl der unterstützten Testausgabeformate bei den meisten untersuchten Testframeworks deutlich steigen. Da keine Aussage über die Vertrauenswürdigkeit und Qualität der externen Plugins getroffen werden kann, beschränkt sich diese Untersuchung auf von \texttt{Nix} bereitgestellte Plugins.
\subsection{Struktur der Testdateien}
Die Struktur der Testdateien wird für jede untersuchte Programmiersprache soweit möglich einheitlich gehalten.\\
Es existieren zwei Testdateien. Sofern das Testframework die Möglichkeit bieten, besitzt jede der Dateien ihren eigenen Namespace. Die Bezeichnung der Namespaces kann variieren, da manche Testframeworks eine bestimmte Namenskonvention voraussetzen.\\
In der ersten Datei liegen 4 Testcases im Namespace \texttt{BasicTest}:
\begin{itemize}
    \item \texttt{test\_neg():} Ein einfacher negativer Testcase, der eine korrekte Assertion beinhaltet und nicht fehlschlägt.
    \item \texttt{test\_pos():} Ein einfacher positiver Testcase, der eine falsche Assertion beinhaltet und fehlschlägt.
    \item \texttt{test\_skip():} Ein Testcase, der übersprungen wird. Für gewönhlich werden übersprungene Testcases trotzdem im der Testausgabe aufgeführt, mit der zusätzlichen Information, dass dieser Testcase übersprungen wurde, anstelle eines Ergebnisses. %Manche Testframeworks unterscheiden zwischen übersprungenen und ignorierten Testcases: Übersprungene Testcases werden nicht ausgeführt, aber als übersprungen in der Testausgabe markiert. Ignorierte Testcases werden weder ausgeführt, noch in die Testausgabe übernommen. Dies macht ignorierte Testcases für die Evaluation der Testausgabeformate uninteressant. 
    \item \texttt{test\_same\_name():} Ein einfacher Testcase, der eine korrekte Assertion beinhaltet und nicht fehlschlägt. In der zweiten Datei existiert ein Testcase mit dem selben Namen.
\end{itemize}
In der zweiten Datei liegt ein Testcase im Namespace \texttt{BasicTest2}:
\begin{itemize}
    \item \texttt{test\_same\_name():} Ein einfacher Testcase, der eine korrekte Assertion beinhaltet und nicht fehlschlägt. In der ersten Datei existiert ein Testcase mit dem selben Namen.
\end{itemize}
\texttt{test\_same\_name} untersucht, ob es allein aus der Testausgabe heraus möglich ist, einen Testcase eindeutig zu identifizieren, auch wenn in unterschiedlichen Namespaces verschiedene Testcases den selben Namen aufweisen.

\tabref{Formatvergleich} zeigt die Ergebnisse der Untersuchung des Datensatzes zur Verbreitung von Testausgabeformaten:
\begin{table}[htbp]
    \centering
\makebox[\textwidth][c]{
    \begin{tabular}{|c|c||c|c|c|c|c|c|}
        \hline
        Framework & Sprache & JUnit XML & \makecell{XML\\(anderes)} & OTR & TAP & JSON & HTML\\
        \hline
        \hline
        catch2 & C++ & \tabYes & \tabYes & \tabNo & \tabYes & \tabNo & \tabNo\\
        \hline
        doctest & C++ & \tabYes & \tabYes & \tabNo & \tabNo & \tabNo & \tabNo\\
        \hline
        gtest & C++ & \tabYes & \tabNo & \tabNo & \tabNo & \tabYes & \tabNo\\
        \hline
        hspec & Haskell & \tabPartial & \tabNo & \tabNo & \tabNo & \tabNo & \tabNo\\
        \hline
        tasty & Haskell & \tabPartial & \tabNo & \tabNo & \tabPartial & \tabNo & \tabNo\\
        \hline
        junit5/jupiter & Java & \tabYes & \tabNo & \tabYes & \tabNo & \tabNo & \tabNo\\
        \hline
        \makecell{junit5/jupiter\\(maven surefire)} & Java & \tabYes & \tabNo & \tabPartial & \tabNo & \tabNo & \tabPartial\\
        \hline
        \makecell{junit5/jupiter\\(gradle test task)} & Java & \tabYes & \tabNo & \tabPartial & \tabNo & \tabNo & \tabYes \\
        \hline
        behat & PHP & \tabYes & \tabNo & \tabNo & \tabNo & \tabYes & \tabNo\\
        \hline
        codeception & PHP & \tabYes & \tabNo & \tabNo & \tabNo & \tabNo & \tabYes\\
        \hline
        phpunit & PHP & \tabYes & \tabNo & \tabYes & \tabNo & \tabNo & \tabYes\\
        \hline
        pytest & Python & \tabYes & \tabNo & \tabNo & \tabPartial & \tabPartial & \tabPartial \\
        \hline
        unittest & Python & \tabYes & \tabNo & \tabNo & \tabNo & \tabNo & \tabNo\\
        \hline
    \end{tabular}
}
    \vspace{1em}
    \begin{tabular}{ll}
        \tabYes & Nativ unterstützt \\
        \tabPartial & Über Plugins unterstützt\\
        \tabNo & Nicht offiziell unterstützt\\
    \end{tabular}
    \caption{Übersicht der Verbreitung bestimmter Formate}
    \label{Formatvergleich}

\end{table}
\newpage

\section{Evaluation der häufigsten Formate}
In diesem Kapitel werden die häufigsten frameworkübergreifend verwendeten Testausgabeformate kurz vorgestellt und nach den in \chapref{Anforderungen Testausgabeformat} festgesetzen Kriterien auf ihre Eignung zur Benutzung als Input für \texttt{testAuditor} evaluiert.

\subsection{JUnit XML}\label{JUNIT}
\subsubsection{Übersicht}
JUnit XML ist ein XML-basiertes Format, welches ursprünglich für das Java-Testframework JUnit entwickelt wurde. Für JUnit XML existiert kein einheitlicher Standard, d.h. es existiert auch kein beschreibendes .xsd Schema, wie sonst für XML-Formate üblich. Trotzdem wurde JUnit XML von vielen Testframeworks in gleicher oder ähnlicher Form adaptiert, wie es ursprünglich für JUnit entworfen wurde. Außerdem existieren viele weitere Schnittstellen zwischen JUnit XML und anderen für das Software-Testing relevante Programme, beispielsweise viele CI/CD-Systeme. %TODO:change "systeme"
Auf diese Weise hat sich JUnit XML zum häufigsten benutze Testausgabeformat für Testframeworks entwickelt.

\subsubsection{Struktur}
Wie bereits erwähnt gibt es bei JUnit XML keine standartisiertes Schema. Allerdings hat sich eine Grundstruktur durchgesetzt, die von nahezu allen Testframeworks, welche JUnit XML als Testausgabeformat anbieten, verwendet wird. Ein typisches Beispiel für diese Grundstruktur ist \lstref{JUNITbehat}:
\begin{lstlisting}[caption={JUnit XML, erzeugt durch das PHP-Testframework behat.\\ ''\textellipsis{}'' bezeichnet gekürzte Stellen.}, label={JUNITbehat}]
<?xml version="1.0" encoding="UTF-8"?>
<testsuites name="default">
    <testsuite name="Basic tests" file= "features/basic.feature" tests="4" skipped="0" failures="2" errors="0" time="0.006">
        <testcase name="test_neg" classname="Basic tests" status="passed" time="0.004" file="features/basic.feature" line="4"></testcase>
    ...
\end{lstlisting}
Das root-element \texttt{<testsuites>} enthält eine Menge von Elementen mit der Bezeichnung \texttt{<testsuite>}, welche jeweils eine Menge von Elementen mit der Bezeichnung \texttt{<testcase>} enthalten. Jedes \texttt{<testcase>} ist korrespondierend zu einem Testcase aus der ursprünglichen Testsuite des zu testenden Programms und enthält für gewöhnlich Informationen über das Testergebnis, den Ablauf, und den ursprünglichen Testcase. Hier ist allerdings wieder wichtig zu erwähnen, dass der Inhalt dieses Elements ebensowenig standartisiert ist, wie der Inhalt der übrigen Elemente. Deutlich wird dieser Sachverhalt in \lstref{JUNITcompare}. Unterschiede in der grundlegenden Struktur zeigen sich z.B. in \lstref{JUNITtasty}. Hier wird das Element \texttt{<testsuite>} verschachtelt verwendet, statt wie sonst üblich nur einmal.

\begin{lstlisting}[caption={direkter Vergleich verschiedener Implementierungen des <testcase>-tags in ausgewählten JUnit-XML Ausgaben. Verglichen wird jeweils die Ausgabe des fehlerlos ausgeführten Testcases \texttt{test\_neg}.\\ ''\textellipsis{}'' bezeichnet gekürzte Stellen.}, label={JUNITcompare}]
(*@\underline{\texttt{junit5 (maven), Java:}}@*)
<testcase name="test_neg" classname="BasicTest" time="0.012"/>
(*@\underline{\texttt{phpunit, PHP:}}@*)
<testcase name="test_neg" file="/home/.../test.php" line="7" class="BasicTest" classname="BasicTest" assertions="1" time="0.000645"/>
(*@\underline{\texttt{gtest, C++:}}@*)
<testcase name="test_neg" file="./test/gtest/test.cpp" line="3" status="run" result="completed" time="0." timestamp="2026-06-29T22:44:17.561" classname="BasicTests" />
\end{lstlisting}

\begin{lstlisting}[caption={JUnit XML, erzeugt durch das Haskell-Testframework tasty mit dem ant-Plugin.\\ ''\textellipsis{}'' bezeichnet gekürzte Stellen.}, label={JUNITtasty}]
<?xml version='1.0'?>
<testsuites errors="0" failures="1" tests="5" time="0.001">
    <testsuite name="AllTests" tests="5">
        <testsuite name="Test1" tests="4">
            <testcase name="test_neg" time="0.000"
                classname="AllTests.Test1" />
        ...
\end{lstlisting}

Aus \lstref{JUNITcompare} ist ersichtlich, dass sich die benutzten Attribute von \texttt{<testcase>} sowie deren Inhalt stark unterscheiden können. Jedoch werden zwei bestimmte Attribute von den meisten Testframeworks gesetzt, und auch in gleicher Weise verwendet:
\begin{itemize}
    \item \texttt{name='' ''} enthält den Namen des Testcase
    \item \texttt{classname='' ''} enthält die Bezeichnung des Namespaces bzw. der Klasse oder Gruppe des Testcase
\end{itemize}
Dieser Trend aus \lstref{JUNITcompare} setzt sich auch in den übrigen im Datensatz untersuchten JUnit XML Ausgaben fort. Es gibt zwar einige wenige Ausnahmen wie exemplarisch in \lstref{JUNITcodeception} dargestellt, diese sind jedoch eher selten. Eine ausführliche Übersicht, wie häufig die eindeutige Zuordnung zwischen Testergebnis und Testcase durch die Attribute \texttt{name} und \texttt{classname} möglich ist, ist in \tabref{Vergleich_Zuordnung} zu finden.

\begin{lstlisting}[caption={JUnit XML für den Testcase \texttt{test\_neg}, erzeugt durch das PHP-Testframework codeception. Codeception benutzt statt dem häufig verwendetem Attribut "classname" das Attribut "class".\\ ''\textellipsis{}'' bezeichnet gekürzte Stellen.}, label={JUNITcodeception}]
<testcase name="test_neg" class="BasicTest" file="/home/.../Test.php" time="0.000252" assertions="1"/>
\end{lstlisting}

\begin{table}[htbp] % [H] to force the table to this position
    \centering
\makebox[\textwidth][c]{
    \begin{tabular}{|c|c||c|c|c|}
        \hline
        Framework & Sprache & \makecell{Zuordnung über\\
        name+classname} & \makecell{Zuordnung über\\alternative Attribute} & \makecell{Namen der\\alternativen Attribute}\\
        \hline
        \hline
        Catch2 & C++ & \tabYes & \tabUnused & \tabUnused\\
        \hline
        doctest & C++ & \tabYes & \tabUnused & \tabUnused\\
        \hline
        gtest & C++ & \tabYes & \tabUnused & \tabUnused\\
        \hline
        hspec & Haskell & \tabYes & \tabUnused & \tabUnused\\
        \hline
        tasty & Haskell & \tabYes & \tabUnused & \tabUnused\\
        \hline
        JUnit5 & Java & \tabYes & \tabUnused & \tabUnused\\
        \hline
        JUnit5 & Java (Maven) & \tabYes & \tabUnused & \tabUnused\\
        \hline
        JUnit5 & Java (Gradle) & \tabYes & \tabUnused & \tabUnused\\
        \hline
        Behat & PHP & \tabYes & \tabUnused & \tabUnused\\
        \hline
        codeception & PHP & \tabNo & \tabYes & name+class\\
        \hline
        PHPUnit & PHP & \tabYes & \tabUnused & \tabUnused\\
        \hline
        pytest & Python & \tabYes & \tabUnused & \tabUnused\\
        \hline
        unittest & Python & \tabYes & \tabUnused & \tabUnused\\
        \hline
    \end{tabular}
}
    \vspace{1em}
    \begin{tabular}{ll}
        \tabYes & Eindeutige Zuordnung möglich \\
        \tabUnused & nicht notwendig\\
        \tabNo & Eindeutige Zuordnung nicht möglich\\
    \end{tabular}
    \caption{Übersicht zur eindeutigen Zuordnung des Testergebnis zum ursprünglichen Testcase über Attribute in JUnit XML}
    \label{Vergleich_Zuordnung}
\end{table}


\subsubsection{Verbreitung}
Da JUnit XML ursprünglich spezifisch für das Testframework Junit entwickelt wurde, folgt die grundelegende Struktur den Gegebenheiten und Anforderungen dieses Frameworks. Zwar folgen viele weitere Testframeworks auch in anderen Sprachen einer ähnlichen Strukur oder können die Junit XML Struktur trotz einiger Unterschiede benutzen, in einzelnen Fällen kommt es jedoch zu Schwierigkeiten und ungewöhnlichen Anpassungen der XML-Struktur.\\
Abgesehen davon ist Junit XML das einzige Testausgabeformat, dass von jedem im Datensatz untersuchten Testframework ausgegeben werden kann. Zumeist sogar nativ, vereinzelt muss die Unterstützung erst über offiziell unterstützte Plugins hinzugefügt werden. % Für einige Testframeworks ist Junit XML sogar das primäre strukturierte Ausgabeformat für ihre Testläufe, zum Beispiel für Junit5. % Den letzten Satz beweisen, oder weglassen
%Bei manchen Testframeworks ist Junit XML sogar die standart-Ausgabe, beispielsweise bei

\subsubsection{Fazit zu JUnit XML}
Aufgrund der weiten Verbreitung und der trotz der fehlenden Standartisierung ähnlichen Umsetzung in vielen Testframeworks, ist JUnit XML ein geeigneter Input für \texttt{testAuditor}, da es in der Praxis trotz einiger potentieller Fehlerquellen mit einer Vielzahl von Testframeworks kompatibel ist.

\subsection{Test Anything Protocol (TAP)}
\subsubsection{Übersicht}
Das Test Anything Protocol wurde ursprünglich für die Programmiersprache Perl entworfen, wurde über die Jahre hinweg jedoch auch für Testframeworks in anderen Sprachen weiterentwickelt und integriert. Anders als bei JUnit XML wird für TAP ein standartisiertes Schema bereitgestellt und gepflegt, siehe \cite{TAP_14}. Anders als die meisten anderen hier untersuchten Testausgabeformate ist TAP keine spezialisierte Struktur eines bereits bestehenden Dateiformates, sondern ein eigenständiges Dateiformat.

\subsubsection{Struktur}
Zu Beginn einer Datei wird die Gesamtzahl aller Testcases der Datei angegeben. Danach folgt eine Auflistung der Testergebnisse mit optionalen Zusatzinformationen, zumeist der Name des Testcase. Übersprungene Tests können über zusätzliche Dekoratoren idikatiert werden, die in TAP als ''Derectives'' bezeichnet werden.\\
Da die meisten Zusatzinformationen optional sind, zeigt \lstref{TAPmin} bereits eine vollständig gültige, minimale TAP-Datei:
\begin{lstlisting}[caption={Minimales TAP-Beispiel}, label={TAPmin}]
    1..1
    ok
\end{lstlisting}
Solch eine Darstellung ist natürlich ungeeignet, da so keine Zuordnung zwischen Testergebnis und Testcase möglich ist. In der Praxis kann dieser Zusammenhang allerdings durch die optionalen Zusatzinformationen gesesetzt werden, wie beispielhaft mit pytest in \lstref{TAPpy} demonstriert:

\begin{lstlisting}[caption={TAP, erzeugt durch pytest.\\ ''\textellipsis{}'' bezeichnet gekürzte Stellen.}, label={TAPpy}]
    # TAP results for test/pytest/test.py
    ok 1 test/pytest/test.py::test_neg
    not ok 2 test/pytest/test.py::test_pos
    ...
    ok 3 test/pytest/test.py::test_skip # SKIP ...
    ok 4 test/pytest/test.py::test_same_name
    1..4
\end{lstlisting}

\subsubsection{Eignung}
Eine eindeutige Zuordnung zwischen Testergebnis und Testcase ist nicht in allen untersuchten Testframeworks, die TAP als Ausgabeformat bieten, möglich. Nehmen wir catch2: Zwei Testscases dürfen in diesem Framework den selben Namen aufweisen, solange sie mit unterschiedlichen Labels getaggt sind. Diese Zusatzinformationen sind allerdings nicht zwingend im TAP-Format erforderlich, und werden vom catch2-Reporter schlicht nicht übernommen. Dies kann eine Ausgabe wie in \lstref{TAPcatch2} zur Folge haben, bei der zwei Testergebnisse mit dem selben Namen bezeichnet werden:
\begin{lstlisting}[caption={TAP, erzeugt durch catch2.\\ ''\textellipsis{}'' bezeichnet gekürzte Stellen.}, label={TAPcatch2}]
    ...
    # test_same_name
    not ok 3 - false
    # test_same_name
    ok 4 - true
    ...
\end{lstlisting}
Dies macht TAP in diesem Falle als Ausgabeformat ungeeignet. Andere Testframeworks wie bespielsweise pytest oder tasty sorgen allerdings für eine eindeutige Zuordnung in ihren jeweiligen Ausgaben. 

\subsubsection{Fazit zu TAP}
TAP ist aufgrund seiner sehr geringen Minimalstruktur und vor allem wegen seiner nur mäßigen Verbreitung eher nicht als Input für \texttt{testAuditor} geeignet. 

\subsection{Open Test Reporting (OTR)}
\subsubsection{Übersicht}
% https://marcphilipp.de/de/startseite/
Open Test Reporting (OTR) wird von dem selben Entwicklerteam, welches auch hinter dem Framework Junit und damit auch hinter JUnit XML steht, als Nachfolger von eben diesem entwickelt. Dabei haben sie laut eigenen Angaben das Ziel, mit OTR die historisch gewachsenen Probleme von JUnit XML (siehe \chapref{JUNIT}) zu umgehen. OTR soll nach Vorstellungen der Entwickler auf lange Sicht JUnit XMLs Platz als häufigstes Ausgabeformat für Testläufe einnehmen. \cite{OTR_milestone}

Anders als die anderen vorstellten Formate gitb es bei OTR zwei mögliche Dateiformate: Es gibt Implementierungen für XML und HTML. Beide Implementierugnen sind im Gegensatz zu JUnit XML standartisiert und folgen einem definiertem Schema. Beim XML-Format wird zwischen zwei Strukture unterschieden: Hierachisch und eventbasiert. Beide Formate nutzen das selbe Set von Kernelementen, unterscheiden sich aber in ihrer Struktur. Für beide Strukturen gibt es Programmiersprachen-spezifische Erweiterungen.\\
Das HTML-Format wird ausgehend von einer vorher erzeugten OTR-XML Datei erzeugt, und ist ebenso erweiterbar. 

\subsubsection{Struktur}
Für diese Untersuchung wird nur die eventbasierte XML-Struktur betrachtet, da diese Struktur das von den im Datenatz untersuchten Testframeworks benutzte Format ist.\\
Im event-basiertem Format werden die Ergebnisse nicht wie in Junit XML oder TAP durch ein einzelnes Element mit optionalen Attributen repräsentiert, sondern über verschiedene Elemente, welche jeweils das Ergebnis eines Event der Testausführung enthalten. Demonstriert ist diese Struktur in \lstref{OTR_example}:

\begin{lstlisting}[caption={OTR von \texttt{test\_neg}, erzeugt durch phpunit.\\ ''\textellipsis{}'' bezeichnet gekürzte Stellen.}, label={OTR_example}]
...
<e:started id="3" parentId="2" name="test_neg" time="2026-06-25T14:11:44.985961Z">
...
<e:finished id="3" time="2026-06-25T14:11:44.986266Z">
...
\end{lstlisting}
Hierbei enthält das Element \texttt{<e:started>} Informationen über den Testcase aus der ursprünglichen Testdatei, während \texttt{<e:finished>} Informationen über das Ergebnis des Testlaufs enthält.\\
Wie auch in JUnit XML und TAP sind die meisten Informationen und Attribute optional, mit einem entscheidenden Unterschied: Aus dem Schema von OTR lässt sich entnehmen, dass die Elemente \texttt{<e:started>} und \texttt{<e:finished>} folgende Attribute aufweisen müssen: \texttt{id}, \texttt{time} und \texttt{name} \cite{OTR_README}. Während \texttt{id} intern in der Ausgabe benötigt wird, um die verschiedenen Events einander zuzuordnen, und \texttt{time} lediglich den Zeitpunkt des Beginns des Events festhält, ist für die Eignung von OTR als Input für \texttt{testAuditor} vor allem \texttt{name} interessant: Dieses Attribut enthält den Namen des Testcase. Jedoch genügt wie bereits bei TAP in \lstref{TAPcatch2} gezeigt das Vorhandensein des bloßen Namens nicht zur eindeutigen Identifizierung eines Testcase. Die dazugehörige Suite wird in OTR innerhalb des Elements \texttt{<metadata>} angegeben. Dieses Element ist jedoch optional. Somit ist die eindeutige Zuordnung von Testcases zu ihrem Ergbnis zwar teilweise in der Struktur festgehalten, aber noch nicht ausreichend garantiert. 

\subsubsection{Eignung}
Da OTR zum Zeitpunkt der Niederschrift dieser Bachelorarbeit noch sehr jung ist und aktiv entwickelt wird, bleibt abzusehen, ob die Entwickler ihr Ziel eines standartiserten, Programiersprachen-agnostischen Testausgebeformates erreichen können.

Außerdem ist die Vebreitung von OTR noch sehr limitiert: Von den in \tabref{Formatvergleich} betrachteten Testframeworks sind lediglich phpunit und junit5 (welches, wie eingangs erwähnt, vom selben Entwicklerteam wie OTR betreut wird) mit jupiter als runner in der Lage, OTR nativ zu erzeugen. Außerdem kann junit5 mit den runnern von maven/gradle zumindest über Plugins eine OTR-Ausgabe erzeugen.

\subsubsection{Fazit zu OTR}
Auch wenn OTR die am besten geeignetste Struktur der bisher untersuchten Testausgabeformate bietet, ist die Verbreitung von OTR zum aktuellen Zeitpunkt noch zu gering, um die framework-agnostik des Verfahrens gewährleisten zu können, sollte es als Input für \texttt{testAuditor} genutzt werden.

% \subsection{Common Test Report Format (CTRF)}
% \subsubsection{Übersicht}

% Ähnlich wie JUnit XML ein spezielles Format für das Dateiformat XML ist, so gibt es auch ein spezielles Testausgabeformat für JSON: Das Common Test Report Format, kurz CTRF. Ähnlich wie OTR ist CTFR zum aktuellen Zeitpunkt ein noch sehr junges Format, welches aktiv weiterentwickelt wird. Dies zeigt sich auch in seiner Verbreitung: Keine der in \tabref{Formatvergleich} untersuchten Testframeworks bieten die Möglichkeit, CTRF zu erzeugen, weder nativ noch über ein offizielles Plugin.\\
% CTRF ist dennoch interresant zu untersuchen, da es zum Einen das einzige größere frameworkagnostische JSON-Format ist, und zum Anderen, da es striktere Vorgaben für seine Struktur als die anderen zuvor untersuchten Testausgabeformate macht.

% \subsubsection{Struktur}
% Was CTRF aus den übrigen Formaten herrausstechen lässt, ist die Tatsache, dass ein Name für Testcases in der Ausgabe zwingend erforderlich ist. 

% \subsubsection{Eignung}

% \subsubsection{Fazit zu CTFR}

\subsection{Andere Formate}
% teamcity, testdox?
Neben den bisher untersuchten ganz oder teilweise strukturierten Formaten bieten viele Testframeworks auch unstrukturierte Ausgaben in unterschiedlichen Dateiformaten an. Konkret wurden bei der Untersuchung des Datensatzes Testausgaben in den folgenden Dateiformaten festgestellt:
\begin{itemize}
    \item XML (ohne JUnit XML oder OTR)
    \item HTML (ohne OTR)
    \item JSON
\end{itemize}
Das Problem bei diesen Testausgaben ist, dass ihre Struktur voll und ganz von dem Testframework abhängig ist, von dem sie erzeugt werden. 

\begin{lstlisting}[caption={Vergleich von der JSON-Ausgabe von pytest und behat.\\ ''\textellipsis{}'' bezeichnet gekürzte Stellen.}, label={JSON_example}]
(*@\underline{\texttt{JSON, von pytest in python erzeugt:}}@*)
{"created": 1782247688.1675658,
"duration": 0.024770736694335938,
"exitcode": 1,
"root": "/home/.../python/src",
"environment": {},
"summary": {
    "passed": 3,
    "failed": 1,
    "skipped": 1,
    "total": 5,
    "collected": 5
},
"collectors": [
...
(*@\underline{\texttt{JSON, von behat in php erzeugt:}}@*)
{"tests": 5,
"skipped": 0,
"failed": 2,
"pending": 0,
"undefined": 0,
"time": 0.011,
"suites": [
...
\end{lstlisting}

Für diese von Grund auf verschiedenen Ausgaben kann offensichtlich kein einheitlicher Parser geschrieben werden.\\
Somit können für die Wahl des Inputs für \texttt{testAuditor} nur die vorangegangen untersuchten Testausgabeformate in Betracht gezogen werden. 


\section{Fazit zur Evaluation der Testausgabenformate}\label{chap:conclusion_testformat_eval}
\tabref{tab:evaluation_testausgabeformate} zeigt eine Übersicht, welches der drei relevanten Testaugabeformate (JUnit XML, TAP, OTR) welche Anforderung aus \chapref{Anforderungen Testausgabeformat} erfüllen kann:
\begin{table}[H]
    \centering
    \makebox[\textwidth][c]{%
    \begin{tabular}{|c||c|c|c|}
        \hline
        Anforderung & JUnit XML & TAP & OTR \\
        \hline
        \hline
        \makecell{ermöglicht Ermittlung aller\\ vorhandenen Testcases}
            & \tabNo
            & \tabNo
            & \tabNo \\
        \hline
        eindeutige Identifizierung von Testcases
            & \tabPartial
            & \tabPartial
            & \tabPartial \\
        \hline
        \makecell{sagt für jeden Testcase einzeln aus,\\ ob er ausgeführt wurde oder nicht}
            & \tabPartial
            & \tabPartial
            & \tabPartial \\
        \hline
        strukturiertes Format
            & \tabPartial
            & \tabYes
            & \tabYes \\
        \hline
        weite Verbreitung
            & \tabYes
            & \tabPartial
            & \tabPartial \\
        \hline
    \end{tabular}%
    }
    \vspace{0.5em}
    \begin{tabular}{ll}
        \tabYes & Erfüllt \\
        \tabPartial & In Teilen erfüllt \\
        \tabNo & Nicht erfüllt \\
    \end{tabular}
    \caption{Evaluation der Testausgabeformate nach den in \chapref{Anforderungen Testausgabeformat} definierten Anforderungen}
    \label{tab:evaluation_testausgabeformate}
\end{table}
% https://xkcd.com/927/\\
Anhand der vorangegangenen Evaluation der unterschiedlichen Testausgabeformate und \tabref{tab:evaluation_testausgabeformate} wird deutlich, dass es keinen eindeutigen Kandidaten für die Wahl als Input für \texttt{testAuditor} gibt, welcher alle in \chapref{Anforderungen Testausgabeformat} gelisteten Anforderungen erfüllt.\\
Ein gemeinsames Problem: Keines der untersuchten Testausgabeformate zeigt die Gesamtheit aller Testcases im Projekt an, sondern auschließlich Testcases, die ausgeführt wurden. Die eindeutige Zuordnung vom Testergebnis zu den Testcases ist in allen drei relevanten Ausgabeformaten nur möglich, wenn die richtigen optionalen Attribute vom jeweiligen Testframework verwednet werden. Da die Gesamtmenger der Testcases nicht bekannt ist, kann nur über die Teilmenge der ausgeführten Testcases eine Aussage getroffen werden, ob sie ausgeführt wurden, während die Menge der nicht ausgeführten Testcases unbekannt bleibt.\\
Unterschiede gibt es bei der Strukturierung der Formate: In diesem Aspekt ist OTR das geeigneste Format, da OTR im Gegensatz zu TAP und JUnit XML das Vorhandensein des Names eines Testcase im Ausgabeformat garantiert. Eine eindeutige Zuordnung kann dadurch allerdings noch nicht garantiert werden, da Zusätzliche Informationen fehlen, konkret die Testsuite oder Klasse, zu der der Testcase gehört.\\
Der wichtigste Aspekt, um einen Testframwork-agnostischen Input für \texttt{testAuditor} zu bestimmen, ist die Verbreitung des Testausgabeformates in verschiedenen Testframeworks. Selbst ein Format, welches alle anderen Anforderungen komplett erfüllt, wäre in diesem Falle nutzlos, wenn es nicht von einer Vielzahl von Testframeworks benutzt wird.\\
Somit fällt die Wahl als Input für \texttt{testAuditor} auf JUnit XML.\\
Trotz aller in \chapref{JUNIT} gezeigten Schwächen wird Junit XML von einer Vielzahl von Testframeworks als Ausgabeformat angeboten. Darüber Hinaus ist die tatsächliche Implementation des Formates von JUnit XML trotz einer fehlenden Standartisierung häufig genug an den kritischen Punkten übereinstimmend, sodass ein Parser geschrieben werden kann, der in einer Vielzahl von Fällen funktioniert.

Natürlich wäre es wünschenswert, ein strukturiertes Format zu verwenden, um eine sichere Aussage über den Anwendungsbereich des Verfahrens treffen zu können. Um die Anforderung der Framework-agnostik so weit wie möglich zu erfüllen, überwiegt hier jedoch das Argument der überragenden Verbreitung von JUnit XML.
